% REST API Reference — editable & compilable online at https://letx.app
% Compiles with pdflatex. A clean REST API reference: endpoint cards with method
% badges, parameter tables, and JSON response samples. Fictional sample API.
\documentclass[11pt]{article}
\usepackage[a4paper,margin=2.2cm]{geometry}
\usepackage[T1]{fontenc}
\usepackage{lato}
\renewcommand{\familydefault}{\sfdefault}
\usepackage{inconsolata}
\usepackage{amsmath}
\usepackage[table]{xcolor}
\usepackage{colortbl}
\usepackage{booktabs}
\usepackage{tabularx}
\usepackage{array}
\usepackage{titlesec}
\usepackage{listings}
\usepackage{tcolorbox}
\tcbuselibrary{skins,breakable}
\usepackage{fancyhdr}
\usepackage[hidelinks]{hyperref}
\setlength{\parindent}{0pt}
\setlength{\parskip}{4pt}

\definecolor{ink}{HTML}{1E2A38}
\definecolor{accent}{HTML}{3B82F6}
\definecolor{getc}{HTML}{16A34A}   % GET green
\definecolor{postc}{HTML}{2563EB}  % POST blue
\definecolor{delc}{HTML}{DC2626}   % DELETE red
\definecolor{light}{HTML}{EEF2F7}
\definecolor{code}{HTML}{F5F7FA}
\definecolor{mist}{HTML}{64748B}

\titleformat{\section}{\Large\bfseries\color{ink}}{}{0em}{}
\titlespacing*{\section}{0pt}{14pt}{6pt}

% method badge: \method{GET}{getc}
\newcommand{\method}[2]{\colorbox{#2}{\color{white}\ttfamily\bfseries\small\,#1\,}}
% endpoint path
\newcommand{\route}[1]{{\ttfamily\normalsize\color{ink} #1}}

\lstset{basicstyle=\ttfamily\footnotesize,backgroundcolor=\color{code},
  frame=single,rulecolor=\color{light},breaklines=true,
  xleftmargin=6pt,framexleftmargin=6pt,columns=fullflexible,
  keywordstyle=\color{accent},showstringspaces=false}

% one endpoint card
\newtcolorbox{apicard}[1]{enhanced,colback=white,colframe=light,
  boxrule=0.8pt,arc=3pt,left=8pt,right=8pt,top=6pt,bottom=6pt,
  title=#1,coltitle=ink,fonttitle=\bfseries,colbacktitle=light}

\pagestyle{fancy}\fancyhf{}
\lhead{\small\color{mist}Orbit API Reference}
\rhead{\small\color{mist}v3.2}
\cfoot{\small\color{mist}Page \thepage\ \textbullet\ \href{https://letx.app}{letx.app}}
\renewcommand{\headrulewidth}{0.3pt}

\emergencystretch=3em\relax
\begin{document}
% ---- header ----
\begin{tcolorbox}[colback=ink,colframe=ink,arc=4pt,left=14pt,top=10pt,bottom=10pt]
{\color{white}\fontsize{20}{22}\selectfont\bfseries Orbit API Reference}\\[3pt]
{\color{accent!60}\small Base URL: \ttfamily https://api.orbit.dev/v3 \quad\textbullet\quad
 Auth: Bearer token \quad\textbullet\quad Format: JSON}
\end{tcolorbox}

\vspace{4pt}
Orbit is a fictional project-management API. All requests require an
\texttt{Authorization: Bearer <token>} header. Responses are JSON; timestamps are ISO-8601.

% ================= ENDPOINT 1 =================
\section{Projects}

\begin{apicard}{\method{GET}{getc}\quad\route{/projects}}
List all projects visible to the authenticated user, with pagination.

\vspace{3pt}\textbf{Query parameters}
\begin{tabularx}{\linewidth}{@{}l l X@{}}
\toprule
\textbf{Name} & \textbf{Type} & \textbf{Description}\\
\midrule
\ttfamily limit  & integer & Max items per page (1--100, default 20).\\
\ttfamily cursor & string  & Opaque pagination cursor from a previous response.\\
\ttfamily status & string  & Filter by \texttt{active}, \texttt{archived}, or \texttt{all}.\\
\bottomrule
\end{tabularx}

\vspace{4pt}\textbf{Response 200}
\begin{lstlisting}
{
  "data": [
    { "id": "prj_8a2f", "name": "Website Redesign",
      "status": "active", "created_at": "2026-06-01T09:12:00Z" }
  ],
  "next_cursor": "eyJvIjoyMH0"
}
\end{lstlisting}
\end{apicard}

\vspace{6pt}
\begin{apicard}{\method{POST}{postc}\quad\route{/projects}}
Create a new project.

\vspace{3pt}\textbf{Body parameters}
\begin{tabularx}{\linewidth}{@{}l l c X@{}}
\toprule
\textbf{Name} & \textbf{Type} & \textbf{Req.} & \textbf{Description}\\
\midrule
\ttfamily name        & string & yes & Human-readable project name.\\
\ttfamily description & string & no  & Optional summary.\\
\ttfamily template\_id & string & no  & Clone from an existing template.\\
\bottomrule
\end{tabularx}

\vspace{4pt}\textbf{Response 201}
\begin{lstlisting}
{ "id": "prj_9c7d", "name": "Mobile App",
  "status": "active", "created_at": "2026-06-28T14:03:00Z" }
\end{lstlisting}
\end{apicard}

% ================= ENDPOINT 2 =================
\section{Tasks}

\begin{apicard}{\method{GET}{getc}\quad\route{/projects/\{id\}/tasks}}
Retrieve tasks for a project.

\vspace{3pt}\textbf{Path parameters}
\begin{tabularx}{\linewidth}{@{}l l X@{}}
\toprule
\textbf{Name} & \textbf{Type} & \textbf{Description}\\
\midrule
\ttfamily id & string & The project identifier (e.g.\ \texttt{prj\_8a2f}).\\
\bottomrule
\end{tabularx}
\end{apicard}

\vspace{6pt}
\begin{apicard}{\method{DELETE}{delc}\quad\route{/tasks/\{id\}}}
Permanently delete a task. This action cannot be undone.

\vspace{4pt}\textbf{Response 204} --- No Content.
\end{apicard}

% ---- status codes ----
\section{Status Codes}
\begin{tabularx}{\linewidth}{@{}l X@{}}
\toprule
\textbf{Code} & \textbf{Meaning}\\
\midrule
\texttt{200 OK}          & Request succeeded.\\
\texttt{201 Created}     & Resource created.\\
\texttt{400 Bad Request} & Invalid parameters.\\
\texttt{401 Unauthorized}& Missing or invalid token.\\
\texttt{404 Not Found}   & Resource does not exist.\\
\texttt{429 Too Many Requests} & Rate limit exceeded; retry after the reset window.\\
\bottomrule
\end{tabularx}
\end{document}
